\section{Antiderivatives and indefinite integrals}

An antiderivative reverses differentiation. If differentiating \(F\) gives \(f\), then integrating \(f\) should recover \(F\), up to a constant.

\begin{definitionbox}
A function \(F\) is an antiderivative of \(f\) on an interval if
\[
F'(x)=f(x)
\]
for every \(x\) in that interval.
\end{definitionbox}

The phrase ``on an interval'' matters. Antiderivatives are usually discussed where the functions involved are defined and behave consistently.

\subsection*{Why the constant appears}

If \(F'(x)=f(x)\), then for every constant \(C\),
\[
(F(x)+C)'=F'(x)+0=f(x).
\]
Thus one antiderivative gives infinitely many antiderivatives. They differ by constants.

This is why the indefinite integral is written as
\[
\int f(x)\,dx=F(x)+C.
\]
The constant \(C\) is not decoration. It records the information lost when differentiating.

\begin{examplebox}
Since
\[
\frac{d}{dx}x^3=3x^2,
\]
we have
\[
\int 3x^2\,dx=x^3+C.
\]
The functions \(x^3\), \(x^3+5\), and \(x^3-17\) all have derivative \(3x^2\).
\end{examplebox}

\subsection*{Checking an integral}

The best way to check an indefinite integral is to differentiate the answer. If the derivative returns the integrand, the antiderivative is correct.

\begin{examplebox}
Claim:
\[
\int (4x^3-2)\,dx=x^4-2x+C.
\]
Check by differentiating:
\[
\frac{d}{dx}(x^4-2x+C)=4x^3-2.
\]
This matches the integrand, so the integral is correct.
\end{examplebox}

This habit is very useful. Integration techniques can be subtle, but differentiation of the final answer is often straightforward.

\subsection*{Indefinite integrals are families}

The expression
\[
\int f(x)\,dx
\]
does not represent one function. It represents the family of all antiderivatives of \(f\). The constant \(C\) describes this family.

\begin{remarkbox}
When the answer to an indefinite integral is written without \(+C\), it is incomplete unless a particular antiderivative has been explicitly requested.
\end{remarkbox}

\begin{summarybox}
When working with indefinite integrals, ask:
\begin{itemize}
\item What function has the given integrand as its derivative?
\item Did I include the constant of integration?
\item Can I check the answer by differentiating?
\item Is the integral asking for one function or a family of functions?
\end{itemize}
\end{summarybox}

Indefinite integration is the systematic search for antiderivatives.